

\begin{enumerate}
	\item 知乎： 关于超螺旋滑模算法的完整证明过程 \url{https://zhuanlan.zhihu.com/p/672355794}
	\item CSDN： 超螺旋滑模控制详细介绍（全网独家） \url{https://blog.csdn.net/qq_41811966/article/details/134869940}
\end{enumerate}

\subsection{稳定性条件概览}

考虑如下的STA误差方程：
\begin{equation}
	\begin{cases}
		\dot{s}=-k_1 \sig{s}^{1/2}+\varphi\\
		\dot{\varphi}=-k_2 \sgn{s}+\dot{\omega}
	\end{cases}
\end{equation}

\begin{table}[!htb]
	\centering
	\begin{tabular}{c|c}
		\hline
		参数条件  & 出处 
		\\\hline    
		$k_1 >0$,
		$k_2>\|\dot{\omega}\|_\infty$& 
		\makecell{	
			引自多体文献\cite{fengFinitetimeDistributedConvex2020}，\\
			可参考单体文献\cite{morenoStrictLyapunovFunctions2012}(en)，\\
			\cite[第五章]{YangGaoJieHuaMoKongZhiLiLunJiQiZaiQianQuDongXiTongZhongDeYingYongYanJiu2016}(zh)
		}
		\\\hline
		$k_1 >\sqrt{6} k_2$,
		$k_2>\|\dot{\omega}\|_\infty$&
		改写自多体文献\cite{chenDistributedFinitetimeDifferentiator2024}
		\\\hline
	\end{tabular}
	\caption{STA收敛条件对比}
\end{table}